\section{Countable laws and effective series}\label{sec:countable}

A countable probability law cannot be inserted into first-order arithmetic as
an actual completed set.  It is presented by primitive-recursive rational
terms and moduli.  The finite partial sums are exact objects of the preceding
sections.

\subsection{Regular sums from absolute-tail moduli}

\begin{definition}[Absolutely summable rational series name]\label{def:abs-series}
An absolute-series name consists of a primitive-recursive rational sequence
\(a_n\) and a primitive-recursive modulus \(N(p)\) such that
\[
 \sum_{n=N(p)}^{M}\abs{a_n}\le2^{-p}
 \qquad\text{for every }M\ge N(p).
\]
The value of the series is represented by the rational intervals
\[
 I_p=\left[\sum_{n<N(p+1)}a_n-2^{-(p+1)},
          \sum_{n<N(p+1)}a_n+2^{-(p+1)}\right].
\]
\end{definition}

The intervals are nested after a harmless monotone closure operation, and
their widths are at most \(2^{-p}\).  Thus they form a regular-real name.  All
interval endpoints and all checks are primitive recursive.

\begin{theorem}[Effective Cauchy convergence]\label{thm:abs-convergence}
Every absolute-series name determines a unique regular real.  Addition,
scalar multiplication, and finite regrouping of such names preserve effective
summability, with moduli computed from the input moduli.
\end{theorem}

\begin{proof}
If \(M,L\ge N(p)\), then
\[
 \abs{\sum_{n<M}a_n-\sum_{n<L}a_n}
 \le\sum_{n=\min(M,L)}^{\max(M,L)-1}\abs{a_n}
 \le2^{-p}.
\]
This proves the Cauchy estimate using only a finite triangle inequality.
For sums of two series, request one additional bit from each tail; for scalar
multiplication by \(c\), request enough extra bits that
\(\abs c\,2^{-q}\le2^{-p}\).  Finite regrouping is exact at every sufficiently
large finite stage.
\end{proof}

\begin{definition}[Countable probability and expectation]\label{def:countable-law}
A countable rational law is a primitive-recursive sequence
\(p_n\ge0\) with an absolute-tail modulus and regular sum one.  A rational
observable \(x_n\) is integrable when \(a_n=p_nx_n\) has an absolute-tail
modulus.  Its expectation is the regular sum
\[
 \E[X]=\sum_{n=0}^{\infty}p_nx_n.
\]
\end{definition}

\begin{theorem}[Effective invariance under reindexing]\label{thm:reindex}
Let \(a_n\) have an absolute-tail modulus.  Suppose
\(\pi:\N\to\N\) is a primitive-recursive bijection equipped with a coverage
modulus \(C(N)\) satisfying
\[
 \{0,\ldots,N-1\}\subseteq
 \{\pi(0),\ldots,\pi(C(N)-1)\}.
\]
Then \(\sum_na_{\pi(n)}\) has the same regular sum as \(\sum_na_n\), with a
primitive-recursively derived modulus.
\end{theorem}

\begin{proof}
Choose \(N=N(p+2)\) from the absolute-tail modulus and
\(K=C(N)\).  The first \(K\) permuted terms contain every original term below
\(N\).  Their difference from the first \(N\) original terms is a finite sum
of terms whose indices are at least \(N\), hence has absolute value at most
\(2^{-(p+2)}\).  Any additional permuted tail is bounded by the same absolute
tail.  The triangle inequality gives the requested \(2^{-p}\) enclosure.
No appeal to a completed permutation set is made; the finite coverage
certificate is the essential constructive data.
\end{proof}

\subsection{Conditional convergence and rearrangement}

\begin{proposition}[Alternating harmonic series]\label{prop:alt-harmonic}
The alternating harmonic partial sums
\[
 A_N=\sum_{n=1}^{N}\frac{(-1)^{n+1}}n
\]
form a regular real with error at most \(1/(N+1)\).  This regular real is the
effectively defined \(\log 2\).
\end{proposition}

\begin{proof}
Group adjacent terms:
\[
 \frac1{2k-1}-\frac1{2k}=\frac1{(2k-1)2k}>0.
\]
The even partial sums increase and the odd partial sums decrease.  Their gap is
the next term, at most \(1/(N+1)\).  Given precision \(p\), bounded search finds
an \(N\) with \(N+1\ge2^p\), producing a regular interval.  We may define
\(\log2\) by this name; equivalence with integral or Taylor definitions is a
separate effective analytic theorem.
\end{proof}

\begin{theorem}[A concrete rearrangement changes the sum]\label{thm:rearrange}
Rearrange the alternating harmonic series in blocks containing two positive
terms followed by one negative term:
\[
 1+\frac13-\frac12+\frac15+\frac17-\frac14+\cdots.
\]
The resulting block partial sums converge effectively to
\(\tfrac32\log2\), not \(\log2\).
\end{theorem}

\begin{proof}
After \(N\) complete blocks, the sum is
\[
 R_N=\sum_{k=1}^{2N}\frac1{2k-1}-\frac12H_N
 =H_{4N}-\frac12H_{2N}-\frac12H_N,
\]
where \(H_n=\sum_{k=1}^n1/k\).  Write
\(H_m=\log m+\gamma+\epsilon_m\), but do not assume an abstract constant
\(\gamma\): use the elementary finite bounds obtained by comparing
\(1/k\) with \(\int_k^{k+1}1/x\,dx\), which show that
\(H_m-\log m\) is a Cauchy name and that the three error terms cancel with an
explicit \(O(1/N)\) bound.  The logarithmic part is
\[
 \log(4N)-\tfrac12\log(2N)-\tfrac12\log N
 =\tfrac32\log2.
\]
Compact rational quadrature for \(1/x\) supplies the needed regular log name
and modulus.  Equivalently, one can compare \(R_N\) directly with three
alternating-harmonic subsequences.  Every comparison is finite at precision
\(p\), with \(N\) chosen larger than a primitive-recursive multiple of
\(2^p\).

The original series is not absolutely summable because its absolute partial
sums are \(H_N\), which exceed any prescribed natural bound by a primitive-
recursive block estimate.  Thus \Cref{thm:reindex} does not apply, exactly as it
should not.
\end{proof}

\subsection{The logarithmic probability example}

\begin{proposition}[A countable law normalized by \(\log2\)]
\label{prop:log-law}
For \(n\ge1\), put
\[
 p_n=\frac{1}{\log2}\frac1{n2^n}.
\]
Then \(\sum_{n\ge1}p_n=1\), and the integer-valued variable \(X(n)=n\) has
\[
 \E[X]=\frac1{\log2}.
\]
Both statements have explicit primitive-recursive tail moduli.
\end{proposition}

\begin{proof}
The finite identity
\[
 \sum_{n=1}^{N}\frac{x^n}{n}
\]
after differentiation is the geometric sum
\(\sum_{n=1}^{N}x^{n-1}=(1-x^N)/(1-x)\).  Integrating its rational polynomial
identity from \(0\) to \(1/2\), with exact polynomial antiderivatives, yields
\[
 \sum_{n=1}^{N}\frac1{n2^n}
 =\int_0^{1/2}\frac{1-t^N}{1-t}\,dt.
\]
The omitted tail is bounded by
\[
 \sum_{n>N}\frac1{n2^n}
 \le\frac1{N+1}\sum_{n>N}2^{-n}
 =\frac{2^{-N}}{N+1}.
\]
The integral limit is the regular \(\log2\).  Hence the \(p_n\) sum to one.
For the expectation,
\[
 \sum_{n=1}^{N}np_n
 =\frac1{\log2}\sum_{n=1}^{N}2^{-n},
\]
and the geometric tail is exactly \(2^{-N}\).  Taking the effective limit gives
\(1/\log2\); the reciprocal is computed by interval bisection after a rational
positive lower bound for \(\log2\) has been certified.
\end{proof}

\subsection{Extended nonnegative expectation}

\begin{definition}[Lower-real expectation]\label{def:lower-expect}
For a nonnegative countable observable, define the increasing rational partial
expectations
\[
 E_N=\sum_{n<N}p_nx_n.
\]
The sequence \((E_N)\) is a lower-real name.  It is finite as a regular real if
an upper-tail modulus is supplied; it is \(+\infty\) when a primitive-recursive
divergence function \(D(B)\) satisfies \(E_{D(B)}\ge B\) for every natural
\(B\).
\end{definition}

\begin{proposition}[St Petersburg game]\label{prop:st-petersburg}
Let the payoff be \(2^n\) with probability \(2^{-n}\) for \(n\ge1\).  The
partial expectation through outcome \(N\) is exactly \(N\); therefore the
extended expectation is \(+\infty\), with divergence function \(D(B)=B\).
\end{proposition}

\begin{proof}
Each term contributes \(2^{-n}2^n=1\).  Thus the first \(N\) terms sum to
\(N\).  This is a quantifier-free induction on \(N\).  The probabilities
normalize because \(\sum_{n=1}^{N}2^{-n}=1-2^{-N}\) with an explicit tail.
\end{proof}

\begin{definition}[Signed extended expectation]\label{def:signed-extended}
For a signed countable observable, form the lower-real expectations of
\(X^+\) and \(X^-\).  If both are finite, define
\(\E[X]=\E[X^+]-\E[X^-]\).  If exactly one is infinite, define the corresponding
signed infinity.  If both are infinite, the expectation is undefined.
\end{definition}

\begin{proposition}[Classification by positive and negative parts]
\label{prop:parts-classification}
The four cases in \Cref{def:signed-extended} are disjoint and exhaustive for a
presentation whose two lower-real names are each certified finite or divergent.
In the finite-finite case, the resulting regular real agrees with the absolute
series \(\sum p_nx_n\), and \(\sum p_n\abs{x_n}\) is finite.
\end{proposition}

\begin{proof}
The certificates make each part decidably classified at the level of the
presentation.  In the finite-finite case,
\(\abs{x}=x^++x^-\) termwise, so the absolute expectation is the sum of two
finite regular reals.  Conversely, an absolute-tail modulus bounds both parts.
The signed identity follows from finite partial sums and passage to regular
limits using their moduli.
\end{proof}

\begin{proposition}[Undefined alternating unbounded example]
\label{prop:undefined-alt}
Let \(p_n=c/n^2\) for \(n\ge1\), where \(c\) is the regular reciprocal of
\(\sum_{n\ge1}1/n^2\), and let \(X(n)=(-1)^{n+1}n\).  Then both
\(\E[X^+]\) and \(\E[X^-]\) diverge, so \(\E[X]\) is undefined even though
symmetric truncations may converge.
\end{proposition}

\begin{proof}
For odd \(n\), the positive contribution is \(c/n\); for even \(n\), the
negative contribution is \(c/n\).  The odd and even harmonic subseries both
diverge.  A primitive-recursive divergence witness follows from grouping
\(2^{r-1}<n\le2^r\): each full block contributes a fixed positive rational
lower bound.  The constant \(c>0\) has a certified rational lower bound from a
finite upper bound on \(\sum1/n^2\).  Thus each part exceeds any prescribed
bound after a computable number of blocks.

The Basel value \(\pi^2/6\) is unnecessary for divergence; only a positive
normalizing regular real is used.  A wholly rational variant takes
\(p_n=1/(n(n+1))\) and the same alternating magnitude \(n\), leading to the
same two divergent parts.
\end{proof}

\subsection{Countable tail sums}

\begin{corollary}[Effective integer tail formula]\label{cor:countable-tail}
Let \(X\ge0\) be integer-valued under a countable law.  Suppose the expectation
tail has a modulus.  Then
\[
 \E[X]=\sum_{k=0}^{\infty}\Pp(X>k)
\]
as equality of regular or lower reals.
\end{corollary}

\begin{proof}
Truncate both the outcome index and the value at finite bounds.  Apply the exact
finite tail formula.  Monotonicity makes the two truncation directions
increasing.  The supplied expectation-tail modulus controls the discarded
weighted values; this also controls the tail-probability double sum because it
is termwise identical at every finite stage.  Hence the regular names agree.
Without a finiteness modulus, the same argument gives equality of lower reals.
\end{proof}
